\clearpage
\section{Experimental results on the DomainBed benchmark}%
\label{sec:domainbed}%%
% \subsection{Benchmark and experimental details}%
% \label{subsec:implem_domainbed}%
\textbf{Datasets.}
We now present our evaluation on DomainBed \cite{gulrajani2021in}.
By imposing the code, the training procedures and the ResNet50 \cite{he51deep} architecture, DomainBed is arguably the fairest benchmark for \ood generalization.
It includes $5$ multi-domain real-world datasets: PACS \cite{li2017deeper}, VLCS \cite{fang2013unbiased}, OfficeHome \cite{venkateswara2017deep}, TerraIncognita \cite{beery2018recognition} and DomainNet \cite{peng2019moment}. \cite{ye2021odbench} showed that \textit{diversity shift dominates in these datasets}. Each domain is successively considered as the target $T$ while other domains are merged into the source $S$. The validation dataset is sampled from $S$, \ie we follow DomainBed's training-domain model selection.
The experimental setup is further described in \Cref{app:domainbed_description}.
Our code is available at \url{https://github.com/alexrame/diwa}.%will be available anonymously at \url{https://anonymous.4open.science/r/DiWA-anonymous-9884/}.%

\textbf{Baselines.}
ERM is the standard Empirical Risk Minimization.
Coral \cite{coral216aaai} is the best approach based on domain invariance.
SWAD (Stochastic Weight Averaging Densely) \cite{cha2021wad} and MA (Moving Average) \cite{arpit2021ensemble} average weights along one training trajectory but differ in their weight selection strategy.
SWAD \cite{cha2021wad} is the current state of the art (SoTA) thanks to it \enquote{overfit-aware} strategy, yet at the cost of three additional hyperparameters (a patient parameter, an overfitting patient parameter and a tolerance rate) tuned per dataset.
In contrast, MA \cite{arpit2021ensemble} is easy to implement as it simply combines all checkpoints uniformly starting from batch $100$ until the end of training.
\reb{Finally, we report the scores obtained in \cite{arpit2021ensemble} for the costly Deep Ensembles (DENS) \cite{Lakshminarayanan2017} (with different initializations)}: we discuss other ensembling strategies in~\Cref{app:wa_vs_ens_analysis}.
% For the sake of fairness, we only report methods with no ;%

\textbf{Our runs.}
ERM and DiWA share the same training protocol in DomainBed: yet, instead of keeping only one run from the grid-search, DiWA leverages $M$ runs.
In practice, we sample $20$ configurations from the hyperparameter distributions detailed in \Cref{tab:hyperparam} and report the mean and standard deviation across $3$ data splits.
For each run, we select the weights of the epoch with the highest validation accuracy.
ERM and MA select the model with highest validation accuracy across the $20$ runs, following standard practice on DomainBed.
\reb{Ensembling (ENS) averages the predictions of all $M=20$ models (with shared initialization).}
DiWA-restricted selects $1 \leq M \leq 20 $ weights with \Cref{alg:pseudo-code} while DiWA-uniform averages all $M=20$ weights.
DiWA$^{\dagger}$ averages uniformly the $M=3\times20=60$ weights from all $3$ data splits.
DiWA$^{\dagger}$ benefits from larger $M$ (without additional inference cost) and from data diversity (see \Cref{app:more_diversity_comparison}).
However, we cannot report standard deviations for DiWA$^{\dagger}$ for computational reasons. Moreover, DiWA$^{\dagger}$ cannot leverage the restricted weight selection, as the validation is not shared across all $60$ weights that have different data splits.

\subsection{Results on DomainBed}
\label{subsec:domain_bed}
We report our \textbf{main results} in \Cref{table:db_all_training}, detailed per domain in \Cref{app:full_results}.
With a randomly initialized classifier, DiWA$^\dagger$-uniform is the best on PACS, VLCS and OfficeHome:
DiWA-uniform is the second best on PACS and OfficeHome.
On TerraIncognita and DomainNet, DiWA is penalized by some bad runs, filtered in DiWA-restricted which improves results on these datasets.
Classifier initialization with linear probing (LP) \cite{kumar2022finetuning} improves all methods on OfficeHome, TerraIncognita and DomainNet.
On these datasets, DiWA$^\dagger$ increases MA by $1.3$, $0.5$ and $1.1$ points respectively.
After averaging, DiWA$^\dagger$ with LP \textit{establishes a new SoTA of $68.0\%$}, improving SWAD by $1.1$ points.%
\begin{table}[b]%
    \vspace{-1.5em}%
    \caption{\textbf{Accuracy ($\%, \uparrow$) on DomainBed with ResNet50} (best in \textbf{bold} and second best \underline{underlined}).}%
    \centering
    \vspace{0.2em}
    \adjustbox{width=\textwidth}{
        \begin{tabular}{llll|cccccc}
            \toprule
             & \textbf{Algorithm}                                & \textbf{Weight selection} & \textbf{Init}                                  & \textbf{PACS}  & \textbf{VLCS}              & \textbf{OfficeHome}        & \textbf{TerraInc}          & \textbf{DomainNet}      & \textbf{Avg}     \\
            \midrule
             & ERM                                               & N/A                       & \multirow{5}{*}{Random}                        & 85.5 $\pm$ 0.2 & 77.5 $\pm$ 0.4             & 66.5 $\pm$ 0.3             & 46.1 $\pm$ 1.8             & 40.9 $\pm$ 0.1          & 63.3             \\
             & Coral \cite{coral216aaai}                         & N/A                       &                                                & 86.2 $\pm$ 0.3 & 78.8 $\pm$ 0.6             & 68.7 $\pm$ 0.3             & 47.6 $\pm$ 1.0             & 41.5 $\pm$ 0.1          & 64.6             \\
             & SWAD \cite{cha2021wad}                            & Overfit-aware             &                                                & 88.1 $\pm$ 0.1 & 79.1 $\pm$ 0.1             & 70.6 $\pm$ 0.2             & 50.0 $\pm$ 0.3             & 46.5 $\pm$ 0.1          & 66.9             \\
             & MA \cite{arpit2021ensemble}                       & Uniform                   &                                                & 87.5 $\pm$ 0.2 & 78.2 $\pm$ 0.2             & 70.6 $\pm$ 0.1             & 50.3 $\pm$ 0.5             & 46.0 $\pm$ 0.1          & 66.5             \\
            & \reb{DENS \cite{Lakshminarayanan2017,arpit2021ensemble}} & \reb{Uniform: $M=6$}            &                                                & \reb{87.6}           & \reb{78.5}                       & \reb{70.8}                       & \reb{49.2}                       & \reb{\textbf{47.7}}           & \reb{66.8}             \\
            \midrule
            \multirow{12}{*}{\begin{turn}{90} Our runs \end{turn}}
             & ERM                                               & N/A                       & \multirow{6}{*}{Random}                        & 85.5 $\pm$ 0.5 & 77.6 $\pm$ 0.2             & 67.4 $\pm$ 0.6             & 48.3 $\pm$ 0.8             & 44.1 $\pm$ 0.1          & 64.6             \\
             & MA \cite{arpit2021ensemble}                       & Uniform                   &                                                & 87.9 $\pm$ 0.1 & 78.4 $\pm$ 0.1             & 70.3 $\pm$ 0.1             & 49.9 $\pm$ 0.2             & 46.4 $\pm$ 0.1          & 66.6             \\
            & \reb{ENS}                                               & \reb{Uniform: $M=20$}           &                                                & \reb{88.0 $\pm$ 0.1} & \reb{78.7 $\pm$ 0.1}             & \reb{70.5 $\pm$ 0.1}             & \reb{51.0 $\pm$ 0.5}             & \reb{47.4 $\pm$ 0.2}          & \reb{67.1}             \\
             & DiWA                                              & Restricted: $M \leq 20$   &                                                & 87.9 $\pm$ 0.2 & \underline{79.2} $\pm$ 0.1 & 70.5 $\pm$ 0.1             & 50.5 $\pm$ 0.5             & 46.7 $\pm$ 0.1          & 67.0             \\
             & DiWA                                              & Uniform: $M=20$           &                                                & 88.8 $\pm$ 0.4 & 79.1 $\pm$ 0.2             & 71.0 $\pm$ 0.1             & 48.9 $\pm$ 0.5             & 46.1 $\pm$ 0.1          & 66.8             \\
             & DiWA$^{\dagger}$                                  & Uniform: $M=60$           &                                                & \textbf{89.0}  & \textbf{79.4}              & 71.6                       & 49.0                       & 46.3                    & 67.1             \\
            \cmidrule{2-10}
             & ERM                                               & N/A                       & \multirow{6}{*}{LP \cite{kumar2022finetuning}} & 85.9 $\pm$ 0.6 & 78.1 $\pm$ 0.5             & 69.4 $\pm$ 0.2             & 50.4 $\pm$ 1.8             & 44.3 $\pm$ 0.2          & 65.6             \\
             & MA \cite{arpit2021ensemble}                       & Uniform                   &                                                & 87.8 $\pm$ 0.3 & 78.5 $\pm$ 0.4             & 71.5 $\pm$ 0.3             & 51.4 $\pm$ 0.6             & 46.6 $\pm$ 0.0          & 67.1             \\
             & \reb{ENS}                                               & \reb{Uniform: $M=20$}           &                                                & \reb{88.1 $\pm$ 0.3} & \reb{78.5 $\pm$ 0.1}             & \reb{71.7 $\pm$ 0.1}             & \reb{50.8 $\pm$ 0.5}             & \reb{47.0 $\pm$ 0.2}          & \reb{67.2}             \\
             & DiWA                                              & Restricted: $M \leq 20$   &                                                & 88.0 $\pm$ 0.3 & 78.5 $\pm$ 0.1             & 71.5 $\pm$ 0.2             & \underline{51.6} $\pm$ 0.9 & \textbf{47.7} $\pm$ 0.1 & 67.5             \\
             & DiWA                                              & Uniform: $M=20$           &                                                & 88.7 $\pm$ 0.2 & 78.4 $\pm$ 0.2             & \underline{72.1} $\pm$ 0.2 & 51.4 $\pm$ 0.6             & 47.4 $\pm$ 0.2          & \underline{67.6} \\
             & DiWA$^{\dagger}$                                  & Uniform: $M=60$           &                                                & \textbf{89.0}  & 78.6                       & \textbf{72.8}              & \textbf{51.9}              & \textbf{47.7}           & \textbf{68.0}    \\
            \bottomrule%
        \end{tabular}%
    }%
    \label{table:db_all_training}%
\end{table}%
%
\clearpage
\begin{wraptable}[9]{hR!}{0.45\textwidth}%
%   \vspace{1em}%
  \caption{\textbf{Accuracy ($\%, \uparrow$) on OfficeHome} domain \enquote{Art} with various objectives.}%
  \centering%
  \adjustbox{width=1.0\linewidth}{%
    \begin{tabular}{lcccc}%
      \toprule
      \textbf{Algorithm} & \textbf{No WA}               & \textbf{MA} & \textbf{DiWA}       & \textbf{DiWA$^{\dagger}$} \\
      \midrule
      ERM                & 62.9 $\pm$ 1.3             & \underline{65.0} $\pm$ 0.2           & 67.3 $\pm$ 0.2             & 67.7                             \\
      Mixup              & \underline{63.1} $\pm$ 0.7 & \textbf{66.2} $\pm$ 0.3              & 67.8 $\pm$ 0.6             & 68.4                             \\
      Coral              & \textbf{64.4} $\pm$ 0.4    & 64.4 $\pm$ 0.4                       & 67.7 $\pm$ 0.2             & 68.2                             \\
      ERM/Mixup          & N/A                        & N/A                                  & 67.9 $\pm$ 0.7             & \underline{68.9}                 \\
      ERM/Coral          & N/A                        & N/A                                  & \underline{68.1} $\pm$ 0.3 & 68.7                             \\
      ERM/Mixup/Coral    & N/A                        & N/A                                  & \textbf{68.4} $\pm$ 0.4    & \textbf{69.1}                    \\
      \bottomrule%
    \end{tabular}}%
  \label{table:db_home0_training}%
\end{wraptable}%%
\textbf{DiWA with different objectives.}
So far we used ERM that does not leverage the domain information.
\Cref{table:db_home0_training} shows that DiWA-uniform benefits from averaging weights trained with Interdomain Mixup \cite{yan2020improve} and Coral \cite{coral216aaai}: accuracy gradually improves as we add more objectives.
Indeed, as highlighted in \Cref{app:more_diversity_comparison}, DiWA benefits from the increased diversity brought by the various objectives.
This suggests a new kind of linear connectivity across models trained with different objectives; the full analysis of this is left for future work.%
\vspace{-0.2em}%
\subsection{Limitations of DiWA}
\label{subsec:limitations}
Despite this success, DiWA has some limitations.
\textit{First}, DiWA cannot benefit from additional diversity that would break the linear connectivity between weights --- as discussed in \Cref{app:wa_vs_ens_analysis}.
\textit{Second}, DiWA (like all WA approaches) can tackle diversity shift but not correlation shift: this property is explained for the first time in \Cref{subsec:analysisbvc} and illustrated in \Cref{app:failure_corr_shift} on ColoredMNIST.%
\vspace{-0.2em}%